\chapter{\python{} dan \Sage{}}

\Sage{} dibangun di atas \python{}, jadi kita akan mulai dari sana.
\python{} adalah bahasa populer\footnote{Dibuat oleh seorang penggemar
Monty Python.} 
dengan gaya sederhana dan pustaka yang andal.
Bahasa ini sering dipakai untuk membuat skrip, sebagai perekat yang menyatukan bagian-bagian terpisah.
% \Sage{} uses it in this way.

Bagian ini mengasumsikan 
bahwa Anda sudah memahami sedikit pemrograman, jadi
di sini Anda hanya akan mempelajari \python{} secukupnya untuk mulai bekerja. 
Untuk pengantar yang lebih menyeluruh, pelajari tutorial \python{} yang sangat baik,
\cite{PythonTeam19b}.



\section{Dasar-dasar Python}
Jalankan \python, misalnya dengan memberikan perintah
\inlinecode{python3}
di terminal.
Beberapa baris informasi awal akan muncul,
diikuti oleh prompt berupa tiga tanda
lebih besar.
\begin{lstlisting}[style=python]
>>>   
\end{lstlisting}
Inilah interpreter \python{}.
Jika Anda mengetik kode 
\python{} lalu menekan \keyboardkey{Enter}, sistem akan
membaca kode tersebut, mengevaluasinya, dan mencetak hasilnya.
Di bawah nanti kita akan melihat cara menulis dan menjalankan program utuh,
tetapi untuk sementara kita akan bereksperimen di dalam interpreter.
Setelah selesai, Anda dapat 
keluar dari prompt ini dengan \textit{\keyboardkey{Ctrl}-D}.

Cobalah memasukkan ekspresi-ekspresi berikut.
\python{} menanggapinya dengan hasil yang ditampilkan
(bintang ganda menyatakan perpangkatan).
\begin{pythonconsole}
2 - (-1)
1 + 2*3
2**3  
\end{pythonconsole}

Salah satu daya tarik \python{} ialah bahwa hal sederhana biasanya tetap sederhana.
Berikut cara mencetak sesuatu ke layar.\footnote{%
  \protect\python{} tersedia dalam dua versi.
  Kita menggunakan versi~3.
  Versi~2 sudah cukup tua, tetapi Anda mungkin masih menjumpainya, misalnya ketika
  mencari kode melalui mesin pencari.
  Salah satu perbedaan keduanya ialah bahwa pada versi~2 pernyataan
  print tidak memakai tanda kurung.}
\begin{pythonconsole}
print(1, "ditambah", 2, "sama dengan", 3)
\end{pythonconsole}
% Often you can debug just by putting in commands to print things, 
% and having a straightforward print operator helps with that. 

Seperti dalam bahasa komputer lain, 
variabel menyediakan tempat bernama untuk menyimpan nilai.
Baris pertama di bawah menaruh $1$ di tempat yang bernama \inlinecode{i},
sedangkan baris kedua mengambil kembali $1$ tersebut.
\begin{pythonconsole}
i = 1
i + 1
\end{pythonconsole}
Dalam sebagian bahasa pemrograman, Anda harus mendeklarasikan `tipe' suatu variabel
sebelum menggunakannya; misalnya, pada contoh di atas Anda harus menyatakan 
bahwa \inlinecode{i} adalah bilangan bulat sebelum dapat menetapkan \inlinecode{i=1}.
Sebaliknya, \python{} menyimpulkan tipe 
dari cara Anda memakai variabel itu\Dash di atas kita menetapkan $1$ kepadanya,
sehingga \python{} menyimpulkan bahwa \inlinecode{i} harus menjadi tempat untuk menyimpan bilangan bulat.
Jika cara pemakaian suatu variabel berubah, \python{} akan 
menyimpulkannya kembali; di sini kita mengubah isi $x$ dari bilangan bulat menjadi string.
\begin{pythonconsole}
x = 1
x
x = 'a'
x
\end{pythonconsole}

Anda dapat melakukan beberapa hal dalam satu baris.
\begin{pythonconsole}
first_day, last_day = 0, 365
first_day
last_day   
\end{pythonconsole}
\python{} menghitung ruas kanan dari kiri ke kanan, lalu menetapkan 
nilai-nilai tersebut kepada variabel di ruas kiri.

Jika Anda melakukan sesuatu yang tidak diperbolehkan, 
\python{} akan mengeluh: sistem menandai masalah dengan 
memunculkan pengecualian dan, pada baris terakhir, menampilkan pesan 
galat.
\begin{pythonconsole}
"a" + 1
\end{pythonconsole}

Selain bilangan bulat, 
\python{} menyediakan bilangan real dan bilangan kompleks.\footnote{%
  Operasi-operasi ini
  dilakukan dengan bilangan titik-mengambang, yaitu sistem 
  yang tertanam dalam perangkat keras komputer Anda untuk
  memodelkan bilangan real.
  Pada contoh ini, perbedaannya tampak
  karena representasi desimal $0.3$ tidak sepenuhnya akurat.
  Walaupun persoalan yang ditimbulkan oleh representasi titik-mengambang sangat menarik\protect\Dash
  lihat \protect\cite{PythonTeam19a} dan \protect\cite{Goldberg91}\protect\Dash
  kita akan mengabaikannya agar perhatian tetap tertuju pada 
  aljabar linear.}
\begin{pythonconsole}
1/5
0.1 + 0.2
(3+2j) - (1-4j)
\end{pythonconsole}
Seperti para insinyur, \python{} menggunakan huruf~$j$ untuk akar
kuadrat dari~$-1$, bukan~$i$ sebagaimana lazimnya dalam matematika.
(Meskipun demikian, \protect\Sage{} mengizinkan Anda memakai~$i$.)

Variabel juga dapat menyatakan nilai kebenaran.
\begin{pythonconsole}
yankees_buruk = True
yankees_buruk
\end{pythonconsole}
Huruf awalnya harus kapital:~bentuknya
\inlinecode{True}
atau \inlinecode{False}, bukan
\inlinecode{true}
atau \inlinecode{false}.
 
Di atas kita melihat beberapa string yang terdiri atas teks di antara tanda petik.
(String-string itu hanya berisi satu karakter, \inlinecode{'a'},
tetapi \python{} tidak mempersoalkan perbedaan antara sebuah karakter
dan string dengan panjang satu.)
Untuk membuat string, Anda dapat memakai tanda petik ganda maupun tunggal, asalkan
tanda pada kedua ujungnya sama. 
Di sini \inlinecode{x} dan \inlinecode{y}
diapit tanda petik ganda, yang dapat dipakai karena masing-masing memuat
tanda petik tunggal, sedangkan string koma dan titik juga diapit tanda petik tunggal. 
\begin{pythonconsole}
x = "Aku berkata, 'Aku Popeye si pelaut'"
y = "Katanya, 'Aku ya aku, dan hanya itulah aku'"
x + ', ' + y + '.'
\end{pythonconsole}
Baris terakhir memperlihatkan bahwa~\inlinecode{+} 
melakukan konkatenasi string.
Baris itu juga memperlihatkan bahwa Anda dapat mencampurkan string bertanda petik tunggal
dengan string bertanda petik ganda.

Untuk membuat baris baru, Anda dapat memakai garis miring-n, \inlinecode{\\n}, 
di dalam string bertanda petik ganda.
Jika Anda membutuhkan banyak baris baru,
masukkan pemisah baris secara langsung ke dalam string yang diapit tiga tanda petik ganda.
\begin{pythonconsole}
a = """JALAN MENUJU KEBIJAKSANAAN
 
Jalan menuju kebijaksanaan?
-- Ya, caranya jelas
dan mudah diucapkan:
Berbuat salah
dan salah
dan salah lagi
tetapi makin sedikit
dan makin sedikit
dan makin sedikit. --Piet Hein"""
\end{pythonconsole} 
Tiga titik yang mengawali baris-baris setelah baris pertama
adalah cara interpreter
\python{} memberi tahu bahwa apa yang Anda ketik belum lengkap:
Anda sudah membuka sebuah string, tetapi belum menutupnya.

\textit{Kamus} \python{} adalah fungsi hingga,
yaitu himpunan hingga
pasangan $\langle\textit{kunci},\textit{nilai}\rangle$,
dengan syarat bahwa tidak ada kunci yang boleh muncul dua kali.
Anda dapat memakai kamus sebagai basis data sederhana.
\begin{pythonconsole}
kata_bilangan = {'satu': 1, 'dua': 2, 'tiga': 3}
kata_bilangan['satu']
kata_bilangan['empat'] = 4  
kata_bilangan
kata_bilangan['dua'] = -1
kata_bilangan
\end{pythonconsole} 
Kamus merupakan bagian penting dalam \python{}, antara lain karena pencarian nilai 
di dalam kamus berlangsung sangat cepat berkat cara data tersebut 
diatur secara internal.
%% But because of that internal organization,
%% if you ask to see a dictionary then \python{} may show its pairs in 
%% a jumbled-up order. % This changes in 3.6 and 3.7 in Cython, but discussion of insertion order seems out of place

Kamus tersusun atas pasangan, sedangkan daftar \python{} adalah urutan entri tunggal.
\begin{pythonconsole}
a = ['alpha', 'beta', 'gamma']
a
b = ['delta']
c = []
\end{pythonconsole}
Daftar dapat memuat apa saja, termasuk daftar lain.
\begin{pythonconsole}
x = 4
a = ['alpha', [True, x]]
a
len(a)
\end{pythonconsole}
Ambil suatu elemen dari daftar dengan menyebutkan indeksnya, yaitu posisinya di dalam daftar,
di antara kurung siku.
Indeks daftar dimulai dari nol; artinya, elemen pertama dalam
daftar bernomor~$0$.
\begin{pythonconsole}
a = ['alpha', 'beta', 'gamma']
a[0]
a[1]
a[0] = 'Alpha'
a[0]
\end{pythonconsole}
Gunakan indeks negatif untuk menghitung dari belakang.
\begin{pythonconsole}
a = ['alpha', 'beta', 'gamma']
a[-1]
\end{pythonconsole}
Menuliskan dua indeks yang dipisahkan titik dua menghasilkan suatu \textit{irisan} 
dari daftar. 
\begin{pythonconsole} 
a = ['alpha', 'beta', 'gamma', 'delta']
a[1:3]
a[1:-1]
a[1:1]
\end{pythonconsole}
Anda dapat memperpanjang sebuah daftar.
\begin{pythonconsole}
c = ['delta']
c.append('epsilon')
c
a = ['beta', 'gamma']
a.insert(0,'alpha')
a+c
\end{pythonconsole}

Sebuah \textit{tupel} menyerupai daftar karena unsur-unsurnya berurutan.
\begin{pythonconsole}
a = ('cak', 'cek', 'cik', 'cuk')
a[0]
\end{pythonconsole}
Namun, berbeda dari daftar, tupel 
tidak dapat dimutasi; isinya tidak dapat berubah.
\begin{pythonconsole}
a = ('cak', 'cek', 'cik', 'cuk')
a[0] = 'ah'
\end{pythonconsole}
Salah satu kegunaan sifat tidak dapat diubah ini ialah bahwa 
tupel dapat menjadi kunci dalam kamus, sedangkan daftar tidak dapat menjadi kunci.
\begin{pythonconsole}
a = ['ke1az', 5418]
b = ('ke1az', 5418)
d = {a: 'aktif'}
d = {b: 'aktif'}
d
\end{pythonconsole}

\python{} mempunyai nilai khusus, \inlinecode{None}, yang
berarti bahwa sesuatu tidak memiliki nilai yang bermakna.
Misalnya, jika program Anda mencatat alamat seseorang dan
memuat variabel \inlinecode{nomor_apartemen}, maka   
untuk orang yang tidak tinggal di
apartemen, Anda sebaiknya menuliskan \inlinecode{nomor_apartemen = None}.



\subsection{Alur kendali}
\python{} mendukung cara-cara tradisional untuk mengubah urutan 
eksekusi pernyataan, dengan sedikit perbedaan.
Perbedaannya ialah bahwa banyak bahasa memakai kurung kurawal atau sintaks lain untuk
menandai suatu blok kode, sedangkan \python{} memakai indentasi.
Selalu gunakan empat spasi untuk mengindentasi. % (Don't use tabs.)
\begin{pythonconsole}
x = 4
if (x == 0):
    delta = 1
else:
    delta = 0

delta
\end{pythonconsole}
Di sini, \python{} menetapkan \inlinecode{delta} menjadi $1$ jika $x$
sama dengan $0$; jika tidak, \python{} menetapkan \inlinecode{delta} menjadi~$0$. 

Setelah \inlinecode{delta = 0}, Anda harus memasukkan 
satu baris kosong agar \python{} mengetahui bahwa blok itu sudah berakhir.
Perhatikan juga bahwa tanda sama dengan ganda, \inlinecode{==}, berarti `sama dengan'. 
% (We cannot use just a single equals because, as we have already seen, 
% \inlinecode{x = 0} 
% means ``$x$ is assigned the value~$0$.'') 
Untuk `tidak sama dengan', gunakan \inlinecode{!=}. 

\python{} memiliki dua variasi lain dari pernyataan \inlinecode{if}.
Pernyataan itu dapat memiliki satu cabang saja,
\begin{pythonconsole}
x = 0
if (x == 0):
    print("Pembagian dengan x tidak diperbolehkan.")

\end{pythonconsole}
atau dapat memiliki lebih dari dua cabang.
\begin{pythonconsole}
x = 2
if (x == 0):
    sgn = 0
elif (x > 0):
    sgn = 1
else:
    sgn = -1

sgn
\end{pythonconsole}

Komputer sangat unggul dalam melakukan iterasi, yaitu mengulang suatu urutan pernyataan.
Perulangan \inlinecode{for} sering melibatkan sebuah \inlinecode{range}.
\begin{pythonconsole}
for i in range(4):
    print(i, "kuadratnya adalah", i**2)

\end{pythonconsole}
Secara bawaan, bilangan terkecil yang dihasilkan oleh \inlinecode{range} adalah $0$,
sesuai dengan penomoran daftar yang dimulai dari nol.
Ubah bilangan awal dengan memanggil \inlinecode{range} 
menggunakan dua argumen.
\begin{pythonconsole}
for i in range(1, 4):
    print(i, "pangkat tiganya adalah", i**3)

\end{pythonconsole}
Bilangan terbesar yang dihasilkan oleh \inlinecode{range} adalah satu kurangnya dari 
masukan terakhirnya, sehingga bilangan terbesar
yang diberikan oleh \inlinecode{range(4)} maupun \inlinecode{range(1,4)} adalah $3$.
Hal ini memudahkan karena gabungan dua daftar
\inlinecode{list(range(4))+list(range(4,10))} menghasilkan 
hal yang sama dengan daftar tunggal \inlinecode{list(range(10))}.

Anda dapat mengiterasi sebuah daftar dengan memakai indeks unsur-unsurnya.
\begin{pythonconsole}
x = [4, 0, 3, 0]
for i in range(len(x)):
    if (x[i] == 0):
        print("butir", i, "bernilai nol")
    else:
        print("butir", i, "tidak bernilai nol")

\end{pythonconsole}
Namun, pengguna \python{} berpengalaman yang tidak sedang mencoba menjelaskan 
perintah \inlinecode{range}
akan mengganti \inlinecode{for i in range(len(x))}
dengan \inlinecode{for c in x}, sebab
\inlinecode{for} dapat mengiterasi urutan apa pun, bukan hanya
daftar bilangan.
Berikut pemformat teks sederhana yang mengisi sebuah baris dengan kata-kata hingga
baris itu akan melebihi $40$ karakter, kemudian
mencetak baris tersebut dan memulai baris baru.
\begin{pythonconsole}
quote = """Pejuang yang menang meraih kemenangan lebih dahulu, lalu pergi berperang, 
sedangkan pejuang yang kalah pergi berperang lebih dahulu, lalu berusaha menang. 
--Sun Tzu"""
line, line_length = [], 0
for wd in quote.split():
    if line_length+len(wd) > 40:
        print(" ".join(line))
        line, line_length = [], 0
    line_length = line_length + len(wd) + 1
    line.append(wd)

if line:
    print(" ".join(line))

\end{pythonconsole}
\python{} memiliki banyak operasi string.
Di atas, \inlinecode{quote.split()} membagi string menjadi kata-kata terpisah, sedangkan
\inlinecode{" ".join(line)} membuat string dengan menyisipkan spasi di antara kata-kata dalam daftar 
\inlinecode{line}.
Dengan demikian, \inlinecode{if line:} terakhir di atas mengeluarkan materi yang masih tersisa
setelah perulangan \inlinecode{for} selesai.

Perulangan \inlinecode{for} dirancang untuk dijalankan sebanyak
jumlah tertentu.
Untuk perulangan yang jumlah pengulangannya belum Anda ketahui sebelumnya, 
gunakan \inlinecode{while}.%% \footnote{%
  %% The Collatz conjecture is 
  %% that this loop will 
  %% terminate for any starting~$n$.
  %% No one knows if the conjecture is true.}
\begin{pythonconsole}
n, i = 12, 0
while (n != 1):
    if (n % 2 == 0):  
        n = n // 2
    else:
        n = 3*n + 1
    i = i + 1
    print("setelah", i, "langkah, n=", n)
    
\end{pythonconsole}
Operasi \inlinecode{//} membagi bilangan bulat di kiri dengan bilangan bulat di kanan, lalu mengembalikan nilai lantainya.

Keluar seketika dari perulangan dengan perintah \inlinecode{break}.
\begin{pythonconsole}
for i in range(10):
    if (i == 3):
        break
    print("di dalam perulangan: i adalah", i)
        
\end{pythonconsole}

Pola perulangan yang umum adalah melakukan suatu tindakan
terhadap setiap unsur daftar.
\python{} memiliki cara ringkas untuk melakukannya, yang disebut komprehensi daftar.
\begin{pythonconsole}
a = [2**n for n in range(4)]
a
[i-1 for i in a]
\end{pythonconsole}



\subsection{Fungsi}
Anda dapat membuat fungsi sendiri.
Berikut implementasi rumus kuadrat.\footnote{%
  Kode ini masih naif; misalnya, kode ini tidak menangani
  potensi persoalan titik-mengambang.}
\begin{pythonconsole}
def quad_formula(a, b, c):
    """Cari akar-akar persamaan kuadrat. 
      a, b, c  bilangan real  Koefisien seperti pada ax^2 + bx + c
    """
    discriminant = (b**2 - 4*a*c)**(0.5)  # pangkat 0.5 adalah akar kuadrat
    r1 = (-1*b-discriminant) / (2.0*a)
    r2 = (-1*b+discriminant) / (2.0*a)
    return (r1, r2)

quad_formula(1, -6, 5)
quad_formula(1, 2, 1)
quad_formula(1, -1, -1)
\end{pythonconsole}
Ketika menulis program dalam \python{}, sebagian besar kode yang Anda tulis
berada di dalam fungsi. 

Fungsi selalu mengembalikan sesuatu; 
jika suatu fungsi tidak pernah mengeksekusi \inlinecode{return}, fungsi itu akan
mengembalikan nilai \inlinecode{None}.
Fungsi juga dapat memuat beberapa pernyataan \inlinecode{return}, misalnya 
satu untuk cabang \inlinecode{if} dan satu untuk cabang \inlinecode{else}.

Mengenai komentar:~fungsi \inlinecode{quad_formula} memiliki dua jenis 
komentar.
Di dalam kode, komentar sebaris diawali tanda pagar, \inlinecode{\#}
(pemrogram sering menulis komentar yang memenuhi satu baris dengan mengawali 
baris itu menggunakan tanda pagar). 
Jenis kedua adalah string bertanda petik tiga yang terletak
setelah baris \inlinecode{def}; string itu
menjelaskan secara singkat
apa yang dilakukan fungsi dan mencantumkan parameternya.

Pada akhir baris \inlinecode{def}, di dalam tanda kurung, terdapat
parameter-parameter fungsi. 
Parameter tersebut menerima nilai 
yang diteruskan pemanggil kepada fungsi.
Fungsi juga dapat mempunyai parameter opsional dengan nilai bawaan.
\begin{pythonconsole}
def sapa(nama="Bapak atau Ibu"):
    """Ucapkan salam.
      nama  string  Nama orang
    """
    print("Halo", nama+".")

sapa("Fred")
sapa()
\end{pythonconsole}
\Sage{} banyak memanfaatkan fitur \python{} ini.







\subsection{Objek dan modul}
Dalam matematika, struktur adalah himpunan objek beserta operasi yang terkait dengannya.
Sebagai contoh, ruang vektor adalah himpunan dengan operasi
penjumlahan dan perkalian skalar.
\python{} berorientasi objek, yang berarti bahwa dengan cara serupa kita dapat menghimpun
data dan tindakan menjadi satu (dalam konteks ini tindakan, yakni fungsi, 
disebut metode).
\begin{pythonconsole}
class DatabaseRecord(object):
    """Simpan informasi tentang seseorang.
    """
    def __init__(self, name, age):
        self.name = name
        self.age = age
    def salutation(self):
        print("Yth.", self.name)

a = DatabaseRecord("Jim", 62)
a.name
a.age
a.salutation()
b = DatabaseRecord("Marie", 20)
b.salutation()
\end{pythonconsole}
Di atas, kita membuat dua instans dari 
\inlinecode{DatabaseRecord}, yaitu instans~\inlinecode{a} dan
instans~\inlinecode{b}.
Untuk memperoleh data usia instans~\inlinecode{a}, tuliskan 
\inlinecode{a.age}.
(Variabel \protect\inlinecode{self} 
dapat membingungkan.
Misalkan pada prompt Anda mengetik \inlinecode{a.name="James"}.
Anda telah memakai nama~\inlinecode{a} 
untuk instans itu, sehingga komputer mengetahui tempat
perubahan tersebut harus dilakukan.
Namun,
di dalam kode deskripsi \inlinecode{class} tidak ada nama instans yang tetap.
Dengan kata lain, \inlinecode{self} merujuk pada instans yang sedang dipakai. 
Sebagai contoh, jika Anda sedang bekerja dengan instans \inlinecode{a}, maka 
di dalam \inlinecode{salutation}, \inlinecode{self.name} merujuk
pada \inlinecode{a.name}.)

Di sini Anda tidak akan menulis kelas sendiri, 
tetapi Anda akan memakai kelas dari
pustaka kode yang ditulis orang lain, termasuk
kode untuk \Sage{}. Oleh karena itu, Anda perlu mengetahui cara memakai 
kelas yang disediakan orang lain.
Dalam konteks ini, pustaka, yaitu kumpulan data, fungsi, dan kelas yang berkaitan,
disebut modul.
Sebagai contoh, \python{} menyertakan modul \inlinecode{math} yang dapat Anda gunakan.
\begin{pythonconsole}
import math
math.pi
math.factorial(5)
math.cos(math.pi)
\end{pythonconsole}
Pernyataan \inlinecode{import} membuat isi
modul tersebut tersedia.

Modul lain menangani bilangan acak.
\begin{pythonconsole}
import random
print(random.randint(1,10))
print(random.randint(0,100)/100)
\end{pythonconsole}



\subsection{Program}
Siklus baca-evaluasi-cetak sangat cocok untuk eksperimen kecil, tetapi
untuk kode yang terdiri atas lebih dari beberapa baris, atau untuk kode yang ingin Anda
simpan, sebaiknya pekerjaan itu ditempatkan dalam berkas terpisah.
Di sini kita akan menguraikan cara menulis program mandiri.

Untuk menulis kode, 
jangan gunakan pengolah kata.
Editor teks menyediakan fitur yang tepat untuk pekerjaan ini.
Gunakan editor yang mendukung \python{}, termasuk
indentasi otomatis dan  
penyorotan sintaks, yaitu pewarnaan kode oleh editor agar lebih mudah
dibaca.
Ada banyak editor yang sangat baik; salah satunya Emacs.\footnote{%
  Program ini mungkin sudah disertakan dalam sistem operasi Anda; jika tidak, 
  lihat \protect\url{www.gnu.org/software/emacs}.}

Berikut contoh pertama program \python{}.
Jalankan editor Anda, buka berkas baru bernama \path{month.py}, lalu masukkan baris-baris berikut.
(Perhatikan bahwa baris pertama menjelaskan program dalam sebuah komentar.
Anda mungkin juga ingin mencantumkan nama penulis, tanggal, dan lisensi kode.
Selain itu, string dokumentasi bertanda petik tiga di bagian atas berkas
merupakan praktik yang baik.
Sebaiknya Anda menyertakan string semacam itu dalam setiap program.)
\begin{lstlisting}[style=python]
# month.py
"""Cetak nomor bulan saat ini. """
import datetime
saat_ini = datetime.datetime.now()  # dapatkan objek datetime, panggil metodenya
print("Nomor bulan saat ini adalah", saat_ini.month)
\end{lstlisting}
Jalankan program itu melalui \python{}.
Salah satu caranya ialah membuka terminal baris perintah,
berpindah ke direktori yang memuat \path{month.py},  
mengetik \inlinecode{python3 month.py}, lalu menekan \keyboardkey{Enter}. 
Anda akan melihat
keluaran seperti \inlinecode{Nomor bulan saat ini adalah 10}.

Berikutnya adalah sebuah permainan kecil. 
(Permainan ini memakai fungsi \python{} \inlinecode{input}, yang menampilkan prompt kepada pengguna
lalu menerima jawaban mereka.)
\lstinputlisting{python/guessing_game.py}
(Fungsi \inlinecode{int} mengubah string dari \inlinecode{input} menjadi bilangan bulat.)
Berikut keluaran ketika permainan ini dijalankan.
\begin{lstlisting}
$ python3 guessing_game.py
Tebak sebuah bilangan bulat antara 1 dan 10: 3
  Maaf, tebakan Anda terlalu kecil
Tebak sebuah bilangan bulat antara 1 dan 10: 8
  Maaf, tebakan Anda terlalu besar
Tebak sebuah bilangan bulat antara 1 dan 10: 5
  Tebakan Anda benar!
\end{lstlisting}  % $

Perhatikan string dokumentasi bertanda petik tiga, baik untuk 
berkas secara keseluruhan maupun untuk fungsi.
Selain berguna bagi pemrogram yang menyunting berkas tersebut, 
string itu terintegrasi dengan sistem bantuan \python{}. 
Masuklah ke direktori yang memuat \path{guessing_game.py}, lalu jalankan 
interpreter \python{}.
Pada prompt \inlinecode{>>>}, masukkan \inlinecode{import guessing_game}.
Anda akan memainkan satu putaran permainan (ada cara untuk mencegah hal ini,
tetapi kita akan mengabaikannya).
Sekarang Anda sedang memakai \inlinecode{guessing_game.py} sebagai modul.
Ketik
\inlinecode{help("guessing_game")}.
Anda akan melihat dokumentasi tentang \inlinecode{guessing_game}, 
termasuk baris-baris berikut
(untuk melihat lebih dari satu halaman, gunakan bilah spasi; untuk keluar dari penampil, tekan \textit{q}).
\begin{lstlisting}
Help on module guessing_game:

NAME
    guessing_game - guessing_game.py

DESCRIPTION
    Tebak bilangan antara 1 dan 10.

FUNCTIONS
    test_guess(guess)
        Tentukan apakah tebakan benar, lalu cetak sebuah pesan.
\end{lstlisting}
\python{} memperoleh informasi ini dari string dokumentasi.
Para pemrogram \Sage{} mengikuti praktik ini, sehingga 
Anda dapat memakai \inlinecode{help(...)} 
kapan pun Anda membutuhkan informasi lebih lanjut tentang rutin \Sage{}.





%----------------------------------
\section{Dasar-dasar \Sage{}}
Sebagian besar manual ini bertujuan memperkenalkan kemampuan \Sage{}.
Namun, karena \Sage{} dibangun di atas \python{}, perbandingan singkat di sini akan berguna.
Lihat \citep{SageTeam19} untuk pengantar yang lebih luas, dan 
ingat pula rujukan \citep{SageTeam19ref}.



\subsection{Baris perintah}
Baris perintah \Sage{} menyerupai baris perintah \python{}, tetapi disesuaikan untuk 
pekerjaan matematika.
Pertama-tama jalankan \Sage{},
misalnya dengan memasukkan \inlinecode{sage} ke jendela baris perintah.
Teks awal dan sebuah prompt akan muncul.
\begin{lstlisting}[style=python]
sage:  
\end{lstlisting}
Keluar dari prompt dengan perintah \inlinecode{exit}
yang diikuti \keyboardkey{Enter}.

Seperti pada Python, Anda dapat memakai interpreter untuk bereksperimen.
\begin{sagecommandline}
sage: 2**3                                                                      
sage: 2^3
sage: 3*1 + 4*2
sage: 5 == 3+3
sage: sin(pi/3)
\end{sagecommandline}
Ekspresi kedua 
memperlihatkan bahwa \Sage{} menambahkan cara ringkas yang praktis untuk perpangkatan, selain
\inlinecode{2**3} milik \python{}.
Ekspresi terakhir
memperlihatkan bahwa \Sage{} kadang-kadang mengembalikan hasil eksak, bukan
hampiran.
Anda tetap dapat memperoleh hampirannya; berikut tiga cara.
\begin{sagecommandline}
sage: sin(pi/3).numerical_approx()
sage: sin(pi/3).n()
sage: n(sin(pi/3), digits=4)  
\end{sagecommandline}
Cara terakhir memperlihatkan adanya opsi untuk membatasi jumlah digit
yang ditampilkan.\footnote{%
  Jika opsi ini tidak digunakan, jumlah digit yang terlihat sesuai dengan 
  presisi $53$ bit, yang merupakan standar untuk  
  bilangan titik-mengambang.}


\subsection{Skrip}
Anda dapat mengelompokkan perintah-perintah \Sage{} dalam sebuah berkas 
agar dapat dipakai kembali tanpa harus mengetiknya ulang.
Buat berkas \path{sage_normal.sage}, 
masukkan fungsi berikut, lalu simpan.
\lstinputlisting{sage/normal_curve.sage}
Jalankan \Sage{}, lalu muat isinya dengan \inlinecode{load(...)}.
\begin{sagecommandline}
sage: load("normal_curve.sage")
sage: normal_curve(1.0)
\end{sagecommandline}


\subsection{Notebook}
Dalam manual ini kita akan tetap menggunakan interpreter.
Namun, perlu disebutkan bahwa \Sage{} juga dapat digunakan melalui   
antarmuka grafis berbasis peramban, 
di dalam Jupyter Notebook, \cite{JupyterTeam19}.
Anda dapat membuat
lembar kerja untuk dipakai sendiri atau bersama orang lain, dengan mudah
kembali dan menyunting perintah sebelumnya, melihat plot yang terintegrasi dengan teks, 
serta memanfaatkan banyak fitur praktis lainnya.
Untuk menjalankannya, alih-alih hanya memberikan perintah \inlinecode{sage},
masukkan
\inlinecode{sage -n jupyter}
(tentu saja, perangkat lunak Jupyter harus sudah
terpasang).\footnote{%
  Jika sistem operasi Anda tidak menyediakannya, kunjungi 
  \protect\url{https://jupyter.org}.}
\endinput
